\scanpage{8}
\mspage{P 51}
intense. L'un et l'autre sont présents en tout moment du travail, et il y a ce mouvement constant de ``va-et\textsuperscript{\unsure{''}}-vient'' entre les ``temps'' où l'un prédomine, et ceux où prédomine l'autre.

\noindent\rule{2.1cm}{0.4pt}

l'autre mode est présent. Quand je construis, aménage, ou que je déblaie, nettoie, ordonne, c'est le ``mode'' ou le ``versant'' ``yang'', ou ``masculin'' du travail qui donne le ton. Quand j'explore à tâtons l'insaisissable, l'informe, ce qui est sans nom, je suis le versant ``yin'', ou ``féminin'' de mon être.

Il n'est pas question pour moi de vouloir minimiser ou renier l'un ou l'autre versant de ma nature, essentiels l'un et l'autre - le ``masculin'' qui construit et qui engendre, et le ``féminin'' qui conçoit, et qui abrite les lentes et obscures gestations. Je ``suis'' l'un et l'autre - ``yang'' et ``yin'', ``homme'' et ``femme''. Mais je sais aussi que l'essence la plus délicate, la plus déliée dans les processus créateurs se trouve du côté du versant ``yin'', ``féminin'' - le versant humble, obscur, et souvent de piètre apparence.

C'est ce versant-là du travail qui, depuis toujours je crois, a exercé sur moi la fascination la plus puissante. Les consensus en vigueur m'encourageaient pourtant à investir le plus clair de mon énergie dans l'autre versant, dans celui qui s'incarne et s'affirme dans des ``produits'' tangibles, pour ne pas dire finis et achevés - des produits aux contours bien tranchés, attestant de leur réalité avec l'évidence de la pierre taillée...

Je vois \qquad bien, avec le recul, comment ces consensus ont pesé sur moi, et aussi comment j'ai ``accusé le poids'' - en souplesse ! La partie ``conception'' ou ``exploration'' de mon travail était maintenue à la portion congrue jusqu'au moment encore de mon départ, soit. Et pourtant, dans ce coup d'oeil rétrospectif sur ce que fut mon oeuvre de mathématicien, il ressort avec une évidence saisissante que ce qui fait l'essence et la puissance de cette oeuvre, c'est bien ce versant de nos jours négligé, quand il n'est objet de dérision ou d'un condescendant dédain : celui des `` i d é e s '' , voire celui du `` r ê v e '' , nullement celui des ``résultats''. Essayant dans ces pages de cerner ce que j'ai apporté de plus essentiel à la mathématique de mon temps, par un regard qui embrasse une forêt, plutôt que de s'attarder sur des arbres - j'ai vu, non un palmarès de ``grands théorèmes'', mais un vivant éventail d'idées fécondes (*), venant concourir toutes à une même et vaste vision.

\noindent\rule{1.8cm}{0.4pt}

(*) Ce n'est pas que ce qu'on peut appeler les ``grands théorèmes'' manquent dans mon oeuvre, y compris des théorèmes qui résolvent des questions posées par d'autres que moi, que personne avant moi n'avait su résoudre. (J'en passe en revue certains dans la note de b. de p. (***) page 554, de la note ``La mer qui monte...'' ( ReS III, n$^\circ$ 122).) Mais, comme je l'ai souligné déjà dès les débuts de
% archive sheet number: 7

